\section{Parabolic and superattracting entropy tails}
\label{sec:tails}

The integer $q=\ell-p$ is the order of the numerator at the projective fixed
point zero.  We now show that $q=1$ and $q\ge2$ produce different certified
remainder laws.

\begin{theorem}[Parabolic boundary]
\label{thm:parabolic}
If $q=1$, then
\begin{equation}
 r_m^\ell=\frac1{\ell m}
 \left\{1+O\left(\frac{\log m}{m}\right)\right\},
 \label{eq:r-parabolic}
\end{equation}
and
\begin{equation}
 h_{\ell,\ell-1}-H_M
 =\frac{\ell^{-M}}{\ell(\ell-1)M}
 \left\{1+O\left(\frac{\log M}{M}\right)\right\}.
 \label{eq:tail-parabolic}
\end{equation}
\end{theorem}

\begin{proof}
Put $s_m=r_m^{-\ell}$.  From
$r_{m+1}=r_m/(1+r_m^\ell)$ we obtain the exact recurrence
\begin{equation}
 s_{m+1}=s_m(1+s_m^{-1})^\ell.
 \label{eq:s-rec}
\end{equation}
Since $s_m\ge s_1=2^\ell$, the binomial theorem gives, uniformly in $m$,
\begin{equation}
 s_{m+1}=s_m+\ell+O(s_m^{-1}).
 \label{eq:s-increment}
\end{equation}
In particular $s_m\ge s_1+\ell(m-1)$, and summing the error in
\cref{eq:s-increment} yields
\begin{equation}
 s_m=\ell m+O(\log m).
 \label{eq:s-asymptotic}
\end{equation}
Taking reciprocals proves \cref{eq:r-parabolic}.

Let $E_M=h_{\ell,\ell-1}-H_M$.  Since
$\log(1+x)=x+O(x^2)$ uniformly for $0\le x\le2^{-\ell}$,
\begin{align}
 E_M
 &=\sum_{k\ge0}\ell^{-(M+k+1)}
 \left\{\frac1{\ell(M+k)}
 +O\left(\frac{\log(M+k)}{(M+k)^2}\right)\right\}.
 \label{eq:E-expand}
\end{align}
For $k\ge0$,
$(M+k)^{-1}=M^{-1}\{1+O(k/M)\}$ whenever $k\le M$, while the part
$k>M$ is exponentially small because of the factor $\ell^{-k}$.
Consequently
\[
 \sum_{k\ge0}\frac{\ell^{-k}}{M+k}
 =\frac1M\left\{\frac{\ell}{\ell-1}+O(M^{-1})\right\}.
\]
The summed error in \cref{eq:E-expand} is
$O(\ell^{-(M+1)}\log M/M^2)$.  Substitution proves
\cref{eq:tail-parabolic}.
\end{proof}

\begin{theorem}[Superattracting regime]
\label{thm:super}
Suppose $q\ge2$ and set $u_m=-\log r_m$.  The positive constant
\begin{equation}
 c_{\ell,p}:=\frac{u_1}{q}
 +\sum_{j=1}^{\infty}q^{-(j+1)}\log(1+r_j^\ell)
 \label{eq:c-constant}
\end{equation}
is finite, and
\begin{equation}
 u_m=c_{\ell,p}q^m+o(1).
 \label{eq:double-exponential}
\end{equation}
Furthermore,
\begin{equation}
 h_{\ell,p}-H_M
 =\ell^{-(M+1)}r_M^\ell\{1+O(r_M^\ell)\}.
 \label{eq:tail-super}
\end{equation}
Thus $r_m$ and the correction beyond its first omitted discounted term are
doubly exponentially small in the level.
\end{theorem}

\begin{proof}
Taking negative logarithms in \cref{eq:r-rec} gives
\begin{equation}
 u_{m+1}=qu_m+\log(1+r_m^\ell).
 \label{eq:u-rec}
\end{equation}
Iteration yields
\begin{equation}
 \frac{u_m}{q^m}=\frac{u_1}{q}
 +\sum_{j=1}^{m-1}q^{-(j+1)}\log(1+r_j^\ell).
 \label{eq:u-sum}
\end{equation}
The summands are positive and bounded by $q^{-(j+1)}2^{-\ell}$, so the
series converges to the positive number in \cref{eq:c-constant}.  The tail
omitted from \cref{eq:u-sum}, after multiplication by $q^m$, is at most
\[
 \sum_{j=m}^{\infty}q^{m-j-1}r_j^\ell.
\]
Because $r_{j+1}\le r_j^q$ and $r_m\to0$, this bound is
$O(r_m^\ell)=o(1)$, proving \cref{eq:double-exponential}.

For the entropy tail, isolate its first term:
\begin{align*}
 h_{\ell,p}-H_M
 &=\ell^{-(M+1)}\log(1+r_M^\ell)
 +\sum_{k\ge1}\ell^{-(M+k+1)}\log(1+r_{M+k}^\ell).
\end{align*}
The first term is
$\ell^{-(M+1)}r_M^\ell\{1+O(r_M^\ell)\}$.
Since $r_{M+k}\le r_M^{q^k}$, the remaining sum is at most
\[
 \ell^{-(M+1)}\sum_{k\ge1}\ell^{-k}r_M^{\ell q^k}
 =\ell^{-(M+1)}r_M^\ell O(r_M^{\ell(q-1)}).
\]
As $q-1\ge1$, this is absorbed by the displayed error.
\end{proof}

\begin{remark}
The discount $\ell^{-m}$ alone is exponential in the level in both regimes.
The distinction is the projective multiplier: polynomial decay at a
parabolic fixed point versus double-exponential decay at a superattracting
one.  This is why the a posteriori bound in \cref{eq:tail-certificate} is
much sharper than a parameter-only geometric bound.
\end{remark}
